% 19.11.2025; 17:11:47


\section[Pattern in der Objekt-Orientierten Programmierung]{OOP Pattern}%
\frame{\sectionpage}
\begin{frame}[fragile]%
    \frametitle{Pattern in der Objekt-Orientierten Programmierung}
    {\large Factory Method Pattern}\vspace{.2em}

    \begin{outline}
        \1[] \begin{minted}[fontsize=\fontsize{7}{8}] {java}
public interface Postage { String type(); }

public class RegularPostage implements Postage {
    public String type() { return "Regular"; }
}

public class ExpressPostage implements Postage {
    public String type() { return "Express"; }
}

public final class PostageFactory {
    private PostageFactory() {}
    public static Postage by_price(double price_for_product) {
        if (price_for_product > 30) {
            return new ExpressPostage();
        } else {
            return new RegularPostage();
        }
    }
}
        \end{minted}
    \end{outline}

\end{frame}
\note {
    \begin{outline}
        \1 Wenn verschiedene Klassen ein \textrm{Interface} implementieren können wir zur Laufzeit entscheiden ein Objekt welcher \textrm{konkreten} Klasse wir zurückgeben
        \1 Methode \texttt{by\_price} von Klasse PostageFactory gibt interface Postage zurück.
        \2 Methode heißt Factory method / ""fabrik methode""
    \end{outline}
}
\begin{frame}
    \begin{outline}
        \1 Wenn verschiedene Klassen ein \textrm{Interface} implementieren können wir zur Laufzeit entscheiden ein Objekt welcher \textrm{konkreten} Klasse wir zurückgeben
        \2 Hier gibt die Methode \texttt{by\_price} der Klasse \texttt{PostageFactory} das interface \texttt{Postage} zurück.
        \2 Diese Methode heißt Factory method oder ""Fabrik Methode""
    \end{outline}
    
\end{frame}


\begin{frame}[fragile]%
    \frametitle{Pattern in der Objekt-Orientierten Programmierung}
    {\large Builder Pattern}\vspace{.2em}

    \begin{outline}
        \1[] \begin{minted}[fontsize=\fontsize{7}{7.4}]{python}
class House:
    def __init__(self, walls, roof, door):
        self.walls = walls
        self.roof = roof
        self.door = door

class HouseBuilder:
    def __init__(self):
        self._walls = "concrete" # Houses have concrete walls by default
        self._roof = "tile"      # Houses have tile roofs by default
        self._door = "glass"     # Houses have glass by default
    def walls(self, style):
        self._walls = style
        return self
    def roof(self, style):
        self._roof = style
        return self
    def door(self, style):
        self._door = style
        return self
    def build(self):
        return House(self._walls, self._roof, self._door)
        \end{minted}
    \end{outline}

\end{frame}
\note {
    \begin{outline}
        \1 Hier Übertragen wir die Verantwortung die Klasse \texttt{House} zu konstruieren an die Klasse \texttt{HouseBuilder}
        \1 Dadurch müssen wir den \texttt{Konstruktor} von \texttt{House} nicht mehr selber  aufrufen und können dadurch Standard Werte festlegen.
        \1 Unnötig in python in diesem kontext
    \end{outline}
}

\begin{frame}[fragile]
    \begin{outline}
        \1 Hier Übertragen wir die Verantwortung die Klasse \texttt{House} zu konstruieren an die Klasse \texttt{HouseBuilder}
        \1 Dadurch müssen wir den \texttt{Konstruktor} von \texttt{House} nicht mehr selber  aufrufen und können dadurch Standard Werte festlegen.
    \end{outline}
    \begin{minted} {python}
       house = HouseBuilder()
                 .walls("brick")
                 .roof("flat")
                 .door("oak")
                 .build()
    \end{minted}
\end{frame}

